\scanpage{7}
\uline{Lemme} Supposons les $L_n$ et $\pi_*(L_{\bullet})$ projectifs de t.f. sur $k$, et supposons $k$ commutatif à \uline{spec connexe} (donc $\forall M$ projectif sur $k$, $\operatorname{rg}_k M\in\mathbb N$ est défini).
% slanted
\margin{on peut \unsure{évid$^{\mathrm{mt}}$} \unsure{éliminer cette|changer cette|éviter cette} hypothèse}

Conditions équivalentes

a) $L_{\bullet}^{!}$ à degrés bornés

b) \deleted{\unsure{$L_{\bullet}$|Tous}} Les $M_i,N_j$ \unsure{n.} nuls sont en nb fini i.e. $L_{\bullet}$ est \inserted{$\simeq$} somme \uline{finie} de modules $\frac12$ simpliciaux $\Lambda^i\Phi_*\otimes M_i$ $(i\geq1)$ $\Lambda^j\Psi_*\otimes N_j$ $(j\geq1)$ ($M_i$ et $N_j$ proj. de t.f.)

b') (Si $k$ principal, plus gén si tout module proj de t.f. sur $k$ est libre). $L_{\bullet}$ est \inserted{$\simeq$} somme \uline{finie} de modules $\frac12$ simpliciaux $\Lambda^i\Phi_*$ \inserted{$(i\geq1)$}, $\Lambda^j\Psi_*$ $(j\geq0)$

c) $\operatorname{rang}_k L_n$ est une fonction polynômiale en $n$ pour $n$ grand

d) $\exists N$ \inserted{et $c>0$} tel que pour $n$ ass. grand, on ait
\[
\operatorname{rang}_k L_n\leq c\,n^N
\]

\deleted{\unsure{Cet}} [De plus, \deleted{p} pour $N$ fixé, la condition que l'on ait $\operatorname{rang}_k L_n=O(n^N)$ signifie que $M_i,N_j$ sont nuls pour $i,j>N$ ; \deleted{et si $N$ est choisi le plus petit possible (i.e. tel que $M_N+N_N\ne0$), \unsure{la limite est un multiple de $1/N!$}.})
% slanted
\margin{et on a plus $\displaystyle\lim_{n\to\infty}\frac{\operatorname{rg}_k L_n}{n^N}=\mu_N+\nu_N$}

\uline{Dém} \deleted{Equivalence de a) et b)} $a)\Longleftrightarrow b)$ déjà vu (voir 1 et 2 ci-dessus), $b)\Longleftrightarrow b')$ trivial, \deleted{$a)\Longrightarrow$} $c)\Longrightarrow d)$ trivial, $a)\Longrightarrow c)$ et $d)\Longrightarrow a)$ immédiat par la formule explicite
% archive sheet number: 6
